\documentclass[11pt]{amsart}

\usepackage[T1]{fontenc}
\usepackage{lmodern}
\usepackage{microtype}
\usepackage{amsmath,amssymb,mathtools}
\usepackage{enumitem}
\usepackage{hyperref}

\hypersetup{
  colorlinks=true,
  linkcolor=blue,
  citecolor=blue,
  urlcolor=blue
}

\newcommand{\N}{\mathbb{N}}
\newcommand{\Q}{\mathbb{Q}}
\newcommand{\Z}{\mathbb{Z}}
\newcommand{\Pui}{\mathcal{P}}
\newcommand{\TrZ}{\operatorname{Tr}_{\Z}}
\newcommand{\den}{\operatorname{den}}
\newcommand{\T}{\mathcal{T}}
\newcommand{\PhiKS}{\Phi_{k,s}}
\newcommand{\qbinom}[2]{\genfrac{[}{]}{0pt}{}{#1}{#2}_{q}}
\newcommand{\coeff}[2]{[q^{#1}]#2}
\newcommand{\geqcoef}{\succeq_q}

\theoremstyle{plain}
\newtheorem{theorem}{Theorem}
\newtheorem{corollary}[theorem]{Corollary}
\newtheorem{proposition}[theorem]{Proposition}
\newtheorem{lemma}[theorem]{Lemma}

\theoremstyle{definition}
\newtheorem{definition}[theorem]{Definition}
\newtheorem{remark}[theorem]{Remark}

\begin{document}



For a nonnegative integer $n$, the $q$-shifted factorial is
\[
  (z;q)_n:=\prod_{j=0}^{n-1}(1-zq^j),
  \qquad (z;q)_0:=1.
\]
For integers $N\geq k\geq 0$, the Gaussian binomial coefficient is
\[
  \qbinom{N}{k}
  :=\frac{(q;q)_N}{(q;q)_k(q;q)_{N-k}}.
\]
It is a polynomial in $q$ with nonnegative integer coefficients.  Among its
standard interpretations, $\qbinom{N}{k}$ is the rank-generating function for
partitions whose Ferrers diagrams fit inside a $k\times (N-k)$ rectangle; see,
for example, \cite[Chapter 3]{Andrews} and \cite[Section 1.7]{Stanley}.

For a positive rational number $r$ and an integer $k\geq 0$, Han and Xiong
\cite{HanXiong} consider the fractional-index extension
\begin{equation}\label{eq:generalized-qbinom}
  G_{r,k}(q)
  :=\qbinom{r+k}{k}
  :=\frac{(q^{r+1};q)_k}{(q;q)_k}.
\end{equation}
When $r$ is a positive integer, this is the usual Gaussian binomial
coefficient.  If $r=a/b$ is nonintegral and written in lowest terms, then
$G_{r,k}(q)$ is naturally a formal power series in $q^{1/b}$.  Han and Xiong
defined its \emph{integer trace} by discarding every term whose exponent is not
an integer.  We write
\[
  \T_{r,k}(q):=\TrZ\bigl(G_{r,k}(q)\bigr)\in\Q[[q]].
\]
For two ordinary power series $F,G\in\Q[[q]]$, write
\[
  F\geqcoef G
  \quad\Longleftrightarrow\quad
  \coeff{n}{(F-G)}\geq 0\quad\text{for every }n\geq 0.
\]
The central conjecture of Han and Xiong is
\begin{equation}\label{eq:HX-conjecture}
  \T_{1/2,k}(q)\geqcoef \T_{r,k}(q)
  \qquad(r\in\Q_{>0},\ k\geq 1).
\end{equation}
They proved \eqref{eq:HX-conjecture} for every $r$ when $k=1,2,3$, and for
every $k$ when $r$ is a positive integer or when $r=m+1/2$ with
$m\geq 0$.  They also verified the cases $r=1/3$ and $r=1/4$ for
$1\leq k\leq 150$ by exact symbolic computation \cite{HanXiong}.

The main result is a general comparison principle.  For a positive rational
number $x$, let $\den(x)$ denote the positive denominator of $x$ in lowest
terms.

\begin{theorem}[Support dominance]\label{thm:support-dominance}
Let $u=c/d$ and $r=a/b$ be positive rational numbers in lowest terms.  Suppose
that
\begin{equation}\label{eq:dominance-hypotheses}
  d\ \text{is even},
  \qquad d\mid \operatorname{lcm}(b,2),
  \qquad u\leq r.
\end{equation}
Then, for every integer $k\geq 1$,
\[
  \T_{u,k}(q)\geqcoef \T_{r,k}(q).
\]
\end{theorem}

Taking $u=1/2$ gives an immediate extension of both infinite families treated
in \cite{HanXiong}.

\begin{corollary}\label{cor:half-line}
For every rational number $r\geq 1/2$ and every integer $k\geq 1$,
\[
  \T_{1/2,k}(q)\geqcoef \T_{r,k}(q).
\]
\end{corollary}

The same theorem gives a reduction of the entire conjecture.  Define the
\emph{even core} of $r$ by
\begin{equation}\label{eq:even-core}
  r^{\flat}:=\frac{1}{\operatorname{lcm}(\den(r),2)}.
\end{equation}
Thus, if $r=a/b$ is reduced, then
\[
 r^{\flat}=
 \begin{cases}
   1/b,& b\text{ even},\\
   1/(2b),& b\text{ odd}.
 \end{cases}
\]

\begin{theorem}[Canonical reduction]\label{thm:canonical-reduction}
For every positive rational number $r$ and every integer $k\geq 1$,
\begin{equation}\label{eq:core-dominates}
  \T_{r^{\flat},k}(q)\geqcoef \T_{r,k}(q).
\end{equation}
Consequently, the Han--Xiong conjecture \eqref{eq:HX-conjecture} is equivalent
to the family of inequalities
\begin{equation}\label{eq:unit-family}
  \T_{1/2,k}(q)\geqcoef \T_{1/(2m),k}(q)
  \qquad(m\geq 1,\ k\geq 1).
\end{equation}
For a fixed $k$, it is enough to check
\begin{equation}\label{eq:finite-unit-family}
  2\leq m\leq \left\lfloor\frac{k}{2}\right\rfloor.
\end{equation}
\end{theorem}



\begin{thebibliography}{99}

\bibitem{Andrews}
G.~E. Andrews,
\emph{The Theory of Partitions},
Cambridge Mathematical Library, Cambridge University Press, Cambridge, 1998.


\bibitem{HanXiong}
G.-N. Han and H.~Xiong,
The $1/2$-conjecture for $q$-binomial coefficients with fractional index,
\emph{arXiv preprint} arXiv:2606.01919v1 (2026).

\bibitem{Stanley}
R.~P. Stanley,
\emph{Enumerative Combinatorics, Vol.~1}, second ed.,
Cambridge Studies in Advanced Mathematics, vol.~49,
Cambridge University Press, Cambridge, 2012.

\end{thebibliography}

\end{document}
